\documentclass[11pt,a4paper]{article}
\usepackage[T1]{fontenc}
\usepackage{lmodern}
\usepackage[margin=2.3cm]{geometry}
\usepackage{amsmath,amssymb}
\usepackage{booktabs}
\usepackage{microtype}
\usepackage{caption}
\usepackage{pgfplots}
\pgfplotsset{compat=1.17}
\usepgfplotslibrary{groupplots}
\usepackage[colorlinks=true,linkcolor=blue!50!black,urlcolor=blue!50!black]{hyperref}
\setlength{\emergencystretch}{3em}

\newcommand{\D}{\mathcal{D}}
\newcommand{\Z}{\mathcal{Z}}
\newcommand{\C}{\mathcal{C}}
\newcommand{\T}{\mathcal{T}}
\newcommand{\pp}{{\uparrow\uparrow}}
\newcommand{\pa}{{\uparrow\downarrow}}
\newcommand{\cin}{c_e^{\rm in}}
\newcommand{\DIR}{\mathrm{DIR}}
\newcommand{\RLC}{\mathrm{RLC}}
\newcommand{\BL}{\mathrm{BL}}

\definecolor{parcol}{HTML}{A32D2D}
\definecolor{antcol}{HTML}{2A78D6}

\title{CQD analysis of the SG0--SG1 spin-resolved beam separations\\[4pt]
\large Data set of 2026-07-06: derivation and fits}
\author{Companion to \emph{Analytical Model and Provisional Calibration of the}\\
\emph{Null Stern--Gerlach CQD Experiment}, Revision F (October 13, 2026)}
\date{October 16, 2026}

\begin{document}
\maketitle

\begin{abstract}
\noindent
The double-stage co-quantum-dynamics (CQD) model of the source note is fitted to the spin-resolved
beam-separation data of 2026-07-06 (\texttt{spin\_resolved\_sep\_20260706.jld2}): 15 SG0 currents
from 0 to 3.64~A for parallel and antiparallel SG0--SG1 fields, each with its own uncertainty. The
chain of equations from the oven ensemble to the detector separation is derived in full, using the
notation of the source note; equation numbers in square brackets refer to it. The SG0 field is
divided by $f=7$ (Revision~F scaling); $f$ is a fixed calibration input, not a fitting parameter.
In the recommended reduction the parallel set alone is fitted with pure CQD, which gives
$k_i=(2.67^{+0.60}_{-0.48})\times10^{-6}$ ($\chi^2=2.8$ for 15 points). With this $k_i$ fixed, pure
CQD predicts the wrong sign for the antiparallel set ($\chi^2=109$); a retained-label fraction
proportional to the SG0 dose, $w_{\rm ret}\simeq0.036~\text{A}^{-1}\times I_0$, reproduces it
($\chi^2=0.6$). With the measured SG0 field ($f=1$) the same procedure gives
$k_i=(1.52\pm0.09)\times10^{-7}$. The alternative fitting strategies are compared, and the
dependence on $f$ is mapped. Fitted curves are supplied as CSV files.
\end{abstract}

\tableofcontents

%==========================================================================
\section{The data}
\label{sec:data}

The beam passes first through SG0, a uniform-field magnet with no gradient, and then through SG1, a
gradient Stern--Gerlach magnet at the fixed current $I_1=40$~mA. The camera records the distance
between the two beam peaks. The data file contains, for SG0 and SG1 fields parallel ($\pp$, group
\texttt{idx\_7}) and antiparallel ($\pa$, group \texttt{idx\_8}), the change of this separation
relative to its zero-current value, $\Delta D(I_0)$ in pixels, together with a standard uncertainty for
every point (Table~\ref{tab:points}, Figure~\ref{fig:data}).

Compared with the digitized eight-point sets used previously, the new data
\begin{itemize}
\item contain 15 currents per configuration, including five low currents $0\le I_0\le0.125$~A and two
near-duplicate pairs (1.403/1.404~A and 3.598/3.642~A);
\item carry individual uncertainties $\sigma_j=0.47$--$0.53$~px, used as weights in every fit;
\item have small zero-current values, $-0.047$~px (parallel) and $+0.049$~px (antiparallel). Their
difference, 0.096~px, equals the stored quantity \texttt{y0}. The model has $\Delta D(0)=0$ by
construction; the offsets are a tenth of $\sigma$ and are kept as data.
\end{itemize}
The model normalizes all separations to the absolute zero-current separation, which is not stored in
the file. The value $D_0=13$~pixels of the source note is used throughout.

\begin{figure}[h]
\centering
\begin{tikzpicture}
\begin{axis}[width=0.86\textwidth,height=0.48\textwidth,xlabel={SG0 current $I_0$ (A)},
  ylabel={$\Delta D$ (px)},xmin=-0.1,xmax=3.8,ymin=-2.8,ymax=5.6,grid=major,grid style={gray!20},
  legend pos=north west,legend style={font=\footnotesize,draw=none,fill=white,fill opacity=0.85,text opacity=1},legend cell align=left]
\addplot[only marks,mark=diamond*,mark size=3pt,parcol,error bars/.cd,y dir=both,y explicit]
  table[x=I0_A,y=dD_parallel_px,y error=sigma_parallel_px,col sep=comma]{spin_resolved_sep_20260706.csv};
\addlegendentry{parallel $\pp$}
\addplot[only marks,mark=square*,mark size=2pt,antcol,error bars/.cd,y dir=both,y explicit]
  table[x=I0_A,y=dD_antiparallel_px,y error=sigma_antiparallel_px,col sep=comma]{spin_resolved_sep_20260706.csv};
\addlegendentry{antiparallel $\pa$}
\addplot[gray,dashed,forget plot] coordinates{(-0.1,0)(3.8,0)};
\end{axis}
\end{tikzpicture}
\caption{The 2026-07-06 data: change of the beam separation relative to zero SG0 current.}
\label{fig:data}
\end{figure}

\paragraph{Physical content.} A uniform field exerts no force, and standard quantum mechanics with the
full hyperfine Hamiltonian predicts no change of the separation (section~23 of the source note). In
CQD, SG0 \emph{rotates the electron magnetic moment} towards one pole without deflecting the atom.
This changes how strongly each atom is pushed inside SG1, and therefore changes the separation. The
size and sign of that change test the CQD collapse law and the electron--nuclear branching rule.

%==========================================================================
